%% kapa: exam=2 round=36 year=2025 period=1 subject=appraisal_practice subject_ko=감정평가실무 seq=4 number=4 points=10
%% source: Q-Net 감정평가사 기출문제 (cst00309) artlSeq=5239325 file=raw/1008/5239325/2025년 제36회 감정평가사 2차 감정평가실무.pdf
%% conversion: kordoc 4.24.51 -> tools/gen2.py (auto); review=proofread-gpt-sol; flags=none
\begin{longquestion}{【문제 4】}
감정평가사 甲은 법원으로 부터 유연분묘가 소재하는 토지에 대한 경매평가를 의뢰받았다. 현장조사 결과 봉분 등이 외부에서 분묘의 존재로 인식될 수 있는 형태임을 확인하였고, 시신의 안장 여부 역시 확인된 경우 다음 물음에 답하시오. (10점) \par
\medskip\noindent 물음 1) 타인의 토지 위에 있는 유연분묘에 “분묘기지권”이 성립될 수 있는 요건을 3가지만 설명하시오. (5점) \par
\medskip\noindent 물음 2) “분묘기지권”이 성립되어 토지에 영향을 미치는 경우 유연분묘가 소재하는 토지의 감정평가 방법을 설명하시오. (5점) \par
\end{longquestion}
